\documentclass[border=12pt]{standalone}
\usepackage{circuitikz}
\usetikzlibrary{arrows.meta}
\begin{document}
\begin{circuitikz}[american, line width=0.9pt]
  % ---- carga trifasica ----
  \draw (6.6,-2.0) rectangle (8.4,2.0);
  \node[align=center] at (7.5,0){carga\\$3\phi$};

  % línea b (referencia, común a las dos bobinas de tensión)
  \draw (0,0) node[left]{$b$} -- (6.6,0);

  % ===================== W1 en la línea a =====================
  \draw (0,1.6) node[left]{$a$} -- (6.6,1.6);
  % flecha del sentido de la corriente que atraviesa la bobina de intensidad
  \draw[-{Latex[length=2.4mm]}] (0.35,2.05) -- (1.05,2.05);
  \node at (0.70,2.35){\small $i_a$};
  % cuerpo del vatímetro (tapa la línea: la línea es la bobina de intensidad, en serie)
  \filldraw[fill=white] (1.55,1.6) circle(0.5);
  \node at (1.55,1.6){$W_1$};
  % borne homólogo (±) de la bobina de intensidad, lado fuente
  \node[circle,fill=black,inner sep=1.0pt] at (1.05,1.6){};
  \node at (0.80,1.92){\small $\pm$};
  % bobina de tensión (en paralelo, referida a b); ± en su línea (a)
  \draw[densely dashed] (1.55,1.1) -- (1.55,0);
  \node[circle,fill=black,inner sep=1.0pt] at (1.55,1.1){};
  \node at (1.82,0.95){\small $\pm$};
  \node[circle,fill=black,inner sep=1.0pt] at (1.55,0){};

  % ===================== W2 en la línea c =====================
  \draw (0,-1.6) node[left]{$c$} -- (6.6,-1.6);
  \draw[-{Latex[length=2.4mm]}] (0.35,-2.05) -- (1.05,-2.05);
  \node at (0.70,-2.35){\small $i_c$};
  \filldraw[fill=white] (1.55,-1.6) circle(0.5);
  \node at (1.55,-1.6){$W_2$};
  \node[circle,fill=black,inner sep=1.0pt] at (1.05,-1.6){};
  \node at (0.80,-1.92){\small $\pm$};
  \draw[densely dashed] (1.55,-1.1) -- (1.55,0);
  \node[circle,fill=black,inner sep=1.0pt] at (1.55,-1.1){};
  \node at (1.82,-0.95){\small $\pm$};
  \node[circle,fill=black,inner sep=1.0pt] at (1.55,0){};

  % leyenda
  \node[align=left,font=\small] at (4.7,2.55){bobina de intensidad: en serie\quad($\pm$ hacia la fuente, por donde entra $i$)};
  \node[align=left,font=\small] at (4.55,-2.75){bobina de tensión: en paralelo a $b$\quad($\pm$ en su misma línea)};
  \node at (4.0,3.15){\small $P_{3\phi}=W_1+W_2\qquad Q_{3\phi}=\sqrt{3}\,(W_2-W_1)$};
\end{circuitikz}
\end{document}
